\subsubsection{Database Model Justification}

The nature of the application's data guided the choice of a relational data model. Core entities such as Users, Expenses, and Reports exhibit clear, well-defined relationships (e.g., a report contains many expenses; an expense belongs to one user). A relational model is particularly effective at enforcing these relationships and ensuring data integrity through mechanisms such as primary keys, foreign keys, and transaction support (ACID compliance). 

\begin{itemize} [leftmargin=0.75cm, noitemsep, label=$\bullet$]

    \item \textbf{Primary keys} define the unique identifier of an entity instance, enabling efficient retrieval of an instance at runtime. 

    \item \textbf{Foreign keys} define relationships between entities, allowing efficient searches based on those relationships when required. 

    \item \textbf{ACID support} guarantees that transactions behave correctly by ensuring the following properties:

    \begin{itemize} [leftmargin=0.75cm, noitemsep, label=$\circ$]
        \item \textbf{Atomicity}: Every transaction is treated as a single unit; either all operations succeed or all fail, preventing partial updates that would produce inconsistent data.

        \item \textbf{Consistency}: Transactions must adhere to the integrity constraints defined for storing valid data.

        \item \textbf{Isolation}: Each transaction is isolated; one transaction does not affect the result of another.

        \item \textbf{Durability}: Once a transaction is successfully committed, its changes are permanently saved, even if the system fails.
    \end{itemize}

\end{itemize}